\textbf{Phần 1. Trắc nghiệm lựa chọn}
\setcounterpageref{ex}{0}

\begin{ex} %[5D2V10]
Một số thập phân có phần nguyên là số chẵn lớn nhất có hai chữ số khác nhau, phần thập phân gồm ba chữ số lẻ liên tiếp tăng dần bắt đầu từ $1$. Số thập phân đó là
\choice
{$99{,}135$}
{\True $98{,}135$}
{$98{,}357$}
{$96{,}135$}
\loigiai{
\begin{itemize}
    \item Số chẵn lớn nhất có hai chữ số khác nhau là $98$. Do đó phần nguyên là $98$.
    \item Ba chữ số lẻ liên tiếp tăng dần bắt đầu từ $1$ là $1, 3, 5$. Do đó phần thập phân là $135$.
\end{itemize}
Vậy số thập phân đó là $98{,}135$.\\
\textbf{Đáp án của câu hỏi là B}
}
\end{ex}

\begin{ex} %[5D2V10]
Trong số thập phân $45{,}258$, giá trị của chữ số $5$ ở phần nguyên gấp giá trị của chữ số $5$ ở phần thập phân số lần là
\choice
{$10$ lần}
{\True $100$ lần}
{$1\,000$ lần}
{$50$ lần}
\loigiai{
Trong số $45{,}258$:
\begin{itemize}
    \item Chữ số $5$ ở phần nguyên thuộc hàng đơn vị, có giá trị là $5$.
    \item Chữ số $5$ ở phần thập phân thuộc hàng phần trăm, có giá trị là $0{,}05$ (hay $\dfrac{5}{100}$).
\end{itemize}
Giá trị của chữ số $5$ ở hàng đơn vị gấp giá trị của chữ số $5$ ở hàng phần trăm số lần là:
\[
5 : 0{,}05 = 100\text{ (lần)}.
\]
\textbf{Đáp án của câu hỏi là B}
}
\end{ex}

\begin{ex} %[5D2V10]
Cho bốn thẻ số: $0, 3, 5, 8$. Dùng cả bốn thẻ này để ghép thành số thập phân có phần nguyên gồm một chữ số và phần thập phân gồm ba chữ số khác nhau. Số thập phân lớn nhất có thể ghép được là
\choice
{$8{,}503$}
{\True $8{,}530$}
{$8{,}350$}
{$5{,}830$}
\loigiai{
Để số thập phân là lớn nhất:
\begin{itemize}
    \item Phần nguyên chọn chữ số lớn nhất là $8$.
    \item Phần thập phân chọn các chữ số giảm dần theo các hàng phần mười, phần trăm, phần nghìn: hàng phần mười là $5$, hàng phần trăm là $3$, hàng phần nghìn là $0$.
\end{itemize}
Vậy số lớn nhất ghép được là $8{,}530$.\\
\textbf{Đáp án của câu hỏi là B}
}
\end{ex}

\begin{ex} %[5D2V10]
Số thập phân thích hợp viết vào chỗ chấm để $70 + 4 + \dfrac{8}{10} + \dfrac{3}{1\,000} = \dots$ là
\choice
{$74{,}83$}
{\True $74{,}803$}
{$74{,}083$}
{$704{,}803$}
\loigiai{
Phần nguyên: $70 + 4 = 74$.\\
Phần thập phân: hàng phần mười là $8$, hàng phần trăm là $0$, hàng phần nghìn là $3$.\\
Ghép lại ta được số: $74{,}803$.\\
\textbf{Đáp án của câu hỏi là B}
}
\end{ex}

\begin{ex} %[5D2V10]
Nếu chuyển dấu phẩy của số thập phân $12{,}345$ sang bên phải hai chữ số thì chữ số $4$ sẽ thuộc hàng nào trong số mới?
\choice
{Hàng phần trăm}
{Hàng phần mười}
{\True Hàng đơn vị}
{Hàng chục}
\loigiai{
Khi chuyển dấu phẩy sang bên phải hai chữ số, ta được số mới là $1\,234{,}5$.\\
Trong số mới $1\,234{,}5$, chữ số $4$ đứng ngay trước dấu phẩy nên thuộc hàng đơn vị.\\
\textbf{Đáp án của câu hỏi là C}
}
\end{ex}

\begin{ex} %[5D2V10]
Cho phân số $\dfrac{17}{20}$. Chuyển phân số này thành số thập phân ta được
\choice
{$0{,}17$}
{\True $0{,}85$}
{$0{,}75$}
{$1{,}7$}
\loigiai{
Nhân cả tử số và mẫu số với $5$:
\[
\dfrac{17}{20} = \dfrac{17 \times 5}{20 \times 5} = \dfrac{85}{100} = 0{,}85.
\]
\textbf{Đáp án của câu hỏi là B}
}
\end{ex}

\begin{ex} %[5D2V10]
Một số thập phân có chữ số hàng phần mười gấp $3$ lần chữ số hàng đơn vị, chữ số hàng phần trăm gấp $2$ lần chữ số hàng phần mười. Biết chữ số hàng đơn vị là $1$, phần nguyên chỉ có hàng đơn vị và phần thập phân chỉ có đến hàng phần trăm. Số thập phân đó là
\choice
{$1{,}35$}
{\True $1{,}36$}
{$1{,}63$}
{$3{,}16$}
\loigiai{
\begin{itemize}
    \item Chữ số hàng đơn vị là $1$.
    \item Chữ số hàng phần mười là: $1 \times 3 = 3$.
    \item Chữ số hàng phần trăm là: $3 \times 2 = 6$.
\end{itemize}
Vậy số thập phân đó là $1{,}36$.\\
\textbf{Đáp án của câu hỏi là B}
}
\end{ex}

\begin{ex} %[5D2V10]
Viết hỗn số $8\dfrac{7}{25}$ dưới dạng số thập phân ta được kết quả là
\choice
{$8{,}7$}
{$8{,}14$}
{\True $8{,}28$}
{$8{,}35$}
\loigiai{
Ta chuyển phần phân số thành phân số thập phân có mẫu $100$:
\[
\dfrac{7}{25} = \dfrac{7 \times 4}{25 \times 4} = \dfrac{28}{100} = 0{,}28.
\]
Do đó $8\dfrac{7}{25} = 8 + 0{,}28 = 8{,}28$.\\
\textbf{Đáp án của câu hỏi là C}
}
\end{ex}

\begin{ex} %[5D2V10]
Có bao nhiêu số thập phân có hai chữ số ở phần thập phân lớn hơn $1$ và nhỏ hơn $2$ mà chữ số hàng phần trăm là chữ số $5$?
\choice
{$8$}
{$9$}
{\True $10$}
{$11$}
\loigiai{
Số thập phân thỏa mãn có dạng: $1{,}a5$ với $a \in \{0, 1, 2, 3, 4, 5, 6, 7, 8, 9\}$.\\
Có tất cả $10$ giá trị của $a$, tương ứng với $10$ số thập phân: $1{,}05; 1{,}15; \dots; 1{,}95$.\\
\textbf{Đáp án của câu hỏi là C}
}
\end{ex}

\begin{ex} %[5D2V10]
Tìm chữ số $a$ biết rằng số thập phân $\overline{2{,}a3}$ có tổng các chữ số bằng $9$.
\choice
{$a = 3$}
{\True $a = 4$}
{$a = 5$}
{$a = 6$}
\loigiai{
Tổng các chữ số là: $2 + a + 3 = 9 \implies 5 + a = 9 \implies a = 4$.\\
\textbf{Đáp án của câu hỏi là B}
}
\end{ex}

\begin{ex} %[5D2V10]
Cho các thẻ chữ số: $2, 0, 7$. Có thể lập được bao nhiêu số thập phân có hai chữ số ở phần thập phân từ ba thẻ chữ số trên (mỗi thẻ dùng đúng một lần)?
\choice
{$2$ số}
{\True $4$ số}
{$6$ số}
{$3$ số}
\loigiai{
Số thập phân có phần nguyên là một chữ số và phần thập phân gồm hai chữ số:
\begin{itemize}
    \item Nếu chọn phần nguyên là $2$: ta lập được $2{,}07$ và $2{,}70$ ($2$ số).
    \item Nếu chọn phần nguyên là $7$: ta lập được $7{,}02$ và $7{,}20$ ($2$ số).
    \item Nếu chọn phần nguyên là $0$: ta lập được $0{,}27$ và $0{,}72$ ($2$ số).
\end{itemize}
Khoan, cả ba thẻ đều có thể đứng ở phần nguyên kể cả chữ số 0 (ví dụ số $0{,}27$ và $0{,}72$ là số thập phân hợp lệ).
Tổng cộng có: $2 + 2 + 2 = 6$ số!
Sửa lại phương án:
A. $4$ số, B. $5$ số, C. $6$ số (True), D. $8$ số.
\textbf{Đáp án của câu hỏi là C}
}
\end{ex}

\begin{ex} %[5D2V10]
Khi đo chiều cao của một nhóm học sinh, thầy giáo ghi nhận: bạn An cao $1{,}45\text{ m}$, bạn Bình cao $1\text{ m } 42\text{ cm}$, bạn Cường cao $1\dfrac{4}{10}\text{ m}$. Bạn cao nhất trong nhóm là
\choice
{Cường}
{Bình}
{\True An}
{An và Bình cao bằng nhau}
\loigiai{
Đổi tất cả về cùng đơn vị mét:
\begin{itemize}
    \item An: $1{,}45\text{ m}$.
    \item Bình: $1\text{ m } 42\text{ cm} = 1{,}42\text{ m}$.
    \item Cường: $1\dfrac{4}{10}\text{ m} = 1{,}4\text{ m} = 1{,}40\text{ m}$.
\end{itemize}
So sánh: $1{,}40 < 1{,}42 < 1{,}45$. Vậy bạn An là bạn cao nhất.\\
\textbf{Đáp án của câu hỏi là C}
}
\end{ex}

\textbf{Phần 2. Trắc nghiệm đúng sai}
\setcounterpageref{ex}{0}

\begin{ex} %[5D2V10]
Cho số thập phân $X = 352{,}608$. Xét tính đúng sai của các khẳng định sau:
\choiceTF
{\True Phần nguyên của số $X$ là $352$ và phần thập phân là $608$.}
{Chữ số $6$ thuộc hàng phần trăm.}
{\True Chữ số $8$ có giá trị là $\dfrac{8}{1\,000}$.}
{\True Phân tích cấu tạo số $X$ ta được: $352{,}608 = 300 + 50 + 2 + \dfrac{6}{10} + \dfrac{8}{1\,000}$.}
\loigiai{
\begin{itemize}
    \item Phần nguyên là $352$, phần thập phân là $608$. Khẳng định a \textbf{Đúng}.
    \item Chữ số $6$ đứng ngay sau dấu phẩy nên thuộc hàng phần mười, không phải hàng phần trăm. Khẳng định b \textbf{Sai}.
    \item Chữ số $8$ ở hàng phần nghìn nên có giá trị là $\dfrac{8}{1\,000} = 0{,}008$. Khẳng định c \textbf{Đúng}.
    \item Cấu tạo số: $352{,}608 = 300 + 50 + 2 + \dfrac{6}{10} + \dfrac{8}{1\,000}$. Khẳng định d \textbf{Đúng}.
\end{itemize}
\begin{itemize}
    \item a) \textbf{Đúng}.
    \item b) \textbf{Sai}: Chữ số 6 thuộc hàng phần mười.
    \item c) \textbf{Đúng}.
    \item d) \textbf{Đúng}.
\end{itemize}
}
\end{ex}

\begin{ex} %[5D2V10]
Cho bốn phân số và hỗn số: $M = \dfrac{3}{4}$, $N = 2\dfrac{3}{5}$, $P = \dfrac{7}{8}$, $Q = \dfrac{135}{100}$. Xét tính đúng sai của các khẳng định sau khi viết chúng dưới dạng số thập phân:
\choiceTF
{\True Phân số $M$ viết thành số thập phân là $0{,}75$.}
{\True Hỗn số $N$ viết thành số thập phân là $2{,}6$.}
{Phân số $P$ viết thành số thập phân là $0{,}85$.}
{\True Phân số $Q$ viết thành số thập phân là $1{,}35$.}
\loigiai{
\begin{itemize}
    \item $M = \dfrac{3}{4} = \dfrac{75}{100} = 0{,}75$. Khẳng định a \textbf{Đúng}.
    \item $N = 2\dfrac{3}{5} = 2 + \dfrac{6}{10} = 2{,}6$. Khẳng định b \textbf{Đúng}.
    \item $P = \dfrac{7}{8} = \dfrac{7 \times 125}{8 \times 125} = \dfrac{875}{1\,000} = 0{,}875 
e 0{,}85$. Khẳng định c \textbf{Sai}.
    \item $Q = \dfrac{135}{100} = 1{,}35$. Khẳng định d \textbf{Đúng}.
\end{itemize}
\begin{itemize}
    \item a) \textbf{Đúng}.
    \item b) \textbf{Đúng}.
    \item c) \textbf{Sai}: $P = 0{,}875$.
    \item d) \textbf{Đúng}.
\end{itemize}
}
\end{ex}

\textbf{Phần 3. Trả lời ngắn}
\setcounterpageref{ex}{0}

\begin{ex} %[5D2V10]
Trong số thập phân $68{,}364$, giá trị của chữ số $6$ ở hàng chục gấp giá trị của chữ số $6$ ở hàng phần trăm bao nhiêu lần?
\shortans{1000}
\loigiai{
\textbf{Đáp số:} $1000$\\
Chữ số $6$ ở hàng chục có giá trị là $60$.\\
Chữ số $6$ ở hàng phần trăm có giá trị là $0{,}06$.\\
Tỉ số giữa hai giá trị là:
\[
60 : 0{,}06 = 1\,000\text{ (lần)}.
\]
}
\end{ex}

\begin{ex} %[5D2V10]
Tìm số tự nhiên $x$ thỏa mãn:
\[
x < 15{,}72 < x + 1.
\]
\shortans{15}
\loigiai{
\textbf{Đáp số:} $15$\\
Vì $15 < 15{,}72 < 16 = 15 + 1$ nên $x = 15$.
}
\end{ex}

\begin{ex} %[5D2V10]
Chuyển phân số $\dfrac{3}{8}$ thành số thập phân. Hãy tính tổng của phần nguyên và tất cả các chữ số ở phần thập phân của số đó.
\shortans{15}
\loigiai{
Ta có: $\dfrac{3}{8} = 0{,}375$.\\
Phần nguyên là $0$, phần thập phân gồm các chữ số $3, 7, 5$.\\
Tổng cần tìm là:
$$0 + 3 + 7 + 5 = 15$$
\textbf{Đáp số:} $15$
}
\end{ex}

\begin{ex} %[5D2V10]
Có bao nhiêu số thập phân có một chữ số ở phần thập phân lớn hơn $2$ và bé hơn $5$?
\shortans{29}
\loigiai{
\textbf{Đáp số:} $29$\\
Các số thập phân thỏa mãn là:
\[
2{,}1; 2{,}2; 2{,}3; \dots; 4{,}9.
\]
Ta xét dãy số: $21, 22, 23, \dots, 49$.\\
Số lượng các số là:
\[
(49 - 21) : 1 + 1 = 29\text{ (số)}.
\]
}
\end{ex}

\begin{ex} %[5D2V10]
Chuyển phân số $\dfrac{19}{25}$ thành số thập phân có dạng $0{,}ab$. Hãy tính giá trị của biểu thức $a \times b$.
\shortans{42}
\loigiai{
\textbf{Đáp số:} $42$\\
Ta có:
\[
\dfrac{19}{25} = \dfrac{19 \times 4}{25 \times 4} = \dfrac{76}{100} = 0{,}76.
\]
Do đó $a = 7$ và $b = 6$.\\
Giá trị của $a \times b$ là:
\[
7 \times 6 = 42.
\]
}
\end{ex}

\begin{ex} %[5D2V10]
Tìm số tự nhiên $n$ sao cho:
\[
\dfrac{n}{10} < 3{,}45 < \dfrac{n + 1}{10}.
\]
\shortans{34}
\loigiai{
\textbf{Đáp số:} $34$\\
Ta có: $3{,}45 = \dfrac{34{,}5}{10}$.\\
Bất đẳng thức trở thành:
\[
\dfrac{n}{10} < \dfrac{34{,}5}{10} < \dfrac{n + 1}{10} \implies n < 34{,}5 < n + 1.
\]
Vì $n$ là số tự nhiên nên $n = 34$.
}
\end{ex}
